% !PRINTED 633
\weberpage{633}
%<<CONT
root of $X = 0$, that is, carries out the substitution $(\varepsilon, \varepsilon^h)$ for $h = 1, 2 \ldots n-1$. All these fields are again contained in $\Omega(\varepsilon)$; they are therefore identical with one another, i.e. $\Omega(\varepsilon)$ is a \gesperrt{normal field}. If a quantity $\Phi(\varepsilon)$ has the property of remaining unchanged when $\varepsilon$ is replaced by all the $\varepsilon^h$, then it is itself contained in $\Omega$. This already follows from the general theorems of \S.\,144. But we see its correctness at once if, under the hypothesis $\Phi(\varepsilon^h) = \Phi(\varepsilon)$, we set
\[
\Phi(\varepsilon) = \frac{1}{n-1}\sum_{1,\,n-1}^{h}\Phi(\varepsilon^h) = c_0 - \frac{c_1 + c_2 + \cdots + c_{n-2}}{n-1}
\]
from which it is then to be concluded, by 2., that $c_1, c_2 \ldots c_{n-2}$ are equal to zero. Hence:

\gesperrt{3. If a quantity in $\Omega(\varepsilon)$ admits all the substitutions $(\varepsilon, \varepsilon^h)$, then it is contained in $\Omega$.}

\begin{center}
\S.\,184.

\gesperrt{Theorems on the Resolvents.}
\end{center}

We now consider the \gesperrt{Lagrange} resolvents $(\varepsilon, x)$ under the hypothesis that $\varepsilon$ is an \gesperrt{imaginary} $n$th root of unity and the $x$ are independent variables, and apply to them the theorems of the preceding section.

We first investigate the influence of the linear permutations on $(\varepsilon, x)$; for this it suffices to determine the influence of the two generating substitutions
\weq{1}{s = (h, h+1), \quad t = (h, gh)}
where $h$, taken modulo $n$, runs through the indices of $x$, and $g$ denotes a primitive root of the prime $n$ (\S.\,180).

But if we apply to
\weq{2}{(\varepsilon, x) = x_0 + \varepsilon x_1 + \varepsilon^2 x_2 + \cdots + \varepsilon^{n-1}x_{n-1}}
the substitution $s$, this function goes over into
\[
x_1 + \varepsilon x_2 + \varepsilon^2 x_3 + \cdots + \varepsilon^{n-1}x_0 = \varepsilon^{-1}(\varepsilon, x).
\]
The effect of the substitution $s$, then, is that
\[
% the German sentence closes with "übergeht." after the display; the verb is absorbed into the
% \text and the display copies the German character for character, without a final stop
(\varepsilon, x) \quad \text{goes over into} \quad \varepsilon^{-1}(\varepsilon, x)
\]
